% element: 10-s059:e04
\BeleskeID{10-s059}
Za određivanje oba ulazna praga polazimo od iste konture:
\[V_I=-V_{CE1}+V_{BE2}+V_{BE3}.\]
Na visokoj granici zamenjujemo režime narednih tranzistora i dobijamo $V_{IH}=-V_{CES}+V_{BE}+V_{BES}$. Na niskoj granici je
\[V_{IL}=-V_{CE1}+V_{BE2}+V_{BE3}=-V_{CES}+V_{\gamma T}+0.\]
Margine za višestruke izvore računaju se neposredno iz $NM_{HMS}=V_{OH}-V_{IH}$ i $NM_{LMS}=V_{IL}-V_{OL}$. Oznake napona i polariteti moraju da se prate duž konkretne putanje, jer visoki izlaz određuje gornji izlazni stepen, a niski izlaz zasićeni $T_3$.
